\MajorSection{القسم الرابع}

\Couplet{٨٤}{أي حقائق التقط ذلك الحكيم الذي استهلكته أقمار كثيرة اكتملت ونقصت؛}{وأي نشيد نبوي له أن ينشده، وماذا كسبت خبرته العتيقة؟}{}

\Couplet{٨٥}{لا إله، لا إله صنعه الإنسان، إنسان أكبر وأقوى وأشد قسوة؛}{شبح أسود من مخاوف طفولتنا، قبل أن يبدأ الفكر، حياة الحياة.}{}

\Couplet{٨٦}{صدق أمير الهند القديم حين قال: «أما أنا فلا أريد إيشوارا؛}{ذلك الخير الأبدي القادر على كل شيء، الذي لا يستطيع أن يخفف الشر الأبدي».}{\BurtonNote{حاشية برتون: بوذا.}}

\Couplet{٨٧}{«قد تكون آلهتكم موجودة، ولكن ما الذي يبيّن أنها موجودة؟» اسمع حكيم الصين الكامل يقول؛}{«وإن كانت موجودة، فما شأنها بنا، وهي تقيم في ذلك الخفاء وذلك البعد؟»}{\BurtonNote{حاشية برتون: كونفوشيوس.}}

\Couplet{٨٨}{«كل مادة لها ميلاد وموت؛ تُصنع وتُنقض وتُصنع من جديد؛}{ونختار أن نسمّي صانعها الله»؛ تلك نظرة الزاهد البومية.}{}

\Couplet{٨٩}{«يا مخلوقات متبدلة متناهية، أنتم تجهدون»؛ يجيب صابّ الخمر؛}{«لتسبروا الأعماق المدوّخة للقدرة اللامتناهية، بمقياسكم المصنوع من خيط».}{\BurtonNote{حاشية برتون: الصوفي أو العارف، في مقابل الزاهد.}}

\Couplet{٩٠}{«وتجهدون لتخلطوا الأصنام الفقيرة، أصنام قلب الإنسان ورأسه، بالفكرة الإلهية؛}{ولتكرزوا بوصفه سير الطبيعة المألوف بما قد يتبدل أو ينتهي في أي ساعة.}{}

\Couplet{٩١}{«كيف يدّعي المعروض أن يعرف شيئًا عن العارض أو العرض؛}{ولماذا تساوم هذا السوم الدنيء لتؤمن، والإيمان لا يعني إلا أنك لن تستطيع المعرفة؟}{}

\Couplet{٩٢}{«كيف يستطيع الآن العابر أن يحتوي الآن الثابت، الأبدية؛}{وجودًا لا نهاية له بلا ماضٍ، وجودًا قائمًا لا يصير أبدًا وجودًا مقبلًا؟}{}

\Couplet{٩٣}{«من صنع صانعكم؟ فإن كان صنع نفسه، فلماذا تذهبون هذا البعد لتبلغوا ما هو أسوأ؛}{ألا يكفي عالم من عوالم، وسلسلة من الكون صنعت نفسها؟}{}

\Couplet{٩٤}{«سلّموا بفكرة، بعلة أولى، بعلة مولّدة للعلل؛ فلماذا تطلبون مزيدًا؛}{ولماذا تجهدون لقياس عمقها وعرضها، وتتبّع عملها، والتماس عونها؟}{}

\Couplet{٩٥}{«مجهول، غير قابل للفهم؛ سمّوه ما تختارون من أسماء؛}{لكن دعوه غائمًا كالفضاء الهوائي، مظلمًا في ظلمته السرية.}{}

\Couplet{٩٦}{«تطلب مخاوفكم الطفولية أبًا، يحدده الإله الذي ليس إنسانًا؛}{تعرفون ما تستطيع حواسكم الخمس معرفته، وتسمّون ما هو كائن تدبيرًا كما يحلو لكم.}{}

\Couplet{٩٧}{«تهبطون بالسماء إلى الأرض المبتذلة، وتصنعون صانعكم على مثال أنفسكم؛}{وترتجفون من الإقرار بسيادة الناموس، وتصلّون إلى الناموس كي ينقض نواميسه.}{}

\Couplet{٩٨}{«تصلّون، ولكن هل وزن فكركم قط ما في الصلاة من فراغ وعبث؛}{وهي تتسول نعمة قد مُنحت بالفعل، أو تطلب أن ترى تبدّلًا في الناموس؟}{}

\Couplet{٩٩}{«قل، أيها الإنسان، يا واسع العلم بالخطة التي تنظم الأسرار السامية؛}{كيف حدث أن كان هذا يسوع، وكان ذاك يهوذا، منذ ميلاد الزمان؟}{}

\Couplet{١٠٠}{«وكيف أكون أنا النمر وتكون أنت الحمل؟ أعد، من فضلك، كشف السر؛}{من قتل القتيل: الرامي أم السهم، أم القدر الذي ساق الرجل والقوس؟}{}

\Couplet{١٠١}{«يعبد الإنسان نفسه؛ وإلهه هو الإنسان؛ إنه كفاح العقل الفاني}{ليصوغ مثاله كما يريده أن يكون، وليجد كمال صورته الخاصة.}{}

\Couplet{١٠٢}{«وصار الإله حكيمًا وكاهنًا وكاتبًا حيث شق ثعبان النيل الوادي؛}{وصار براهم الكئيب في الهند المتوهجة، شيئًا محايدًا باردًا شاحبًا.}{}

\Couplet{١٠٣}{«وصار في جبال كلدان العالية صانعًا للأفلاك السماوية؛}{وفي سهوب المجوس إلهًا خارج الزمان، يحكم بواسطة قرينيه الاثنين.}{}

\Couplet{١٠٤}{«وصار في خيام العبرانيين الرب الذي قاد عبيده البرص إلى القتال والخصام؛}{يهوه أو أدون أو إلوهيم، الإله الذي يضرب، رجل الحرب.}{\BurtonNote{حاشية برتون: يهوه هو \textenglish{Jehovah}.}}

\Couplet{١٠٥}{«أما آلهة اليونان الجميلة المنطلقة من القيود، تلك الصور الإنسانية الحسنة الهشّة؛}{فقد كانت مساكنها تطل على البحر الأوسط، حيث يرقد جمال الأرض كله في مهده.}{}

\Couplet{١٠٦}{«ولم تغادر قط حدوده المباركة، ولم تطلب أقاليم همجية لآلهة همجية؛}{حيث يهز أودين الشمال الكئيب رأسه فوق الخنزير وكأس شراب العسل المَقيت.}{}

\Couplet{١٠٧}{«وحين ارتفعت أخيرًا الصيحة العالية الموجعة: مات بان العظيم؛}{ذبل سحر على الأرض، وخبا بهاء في السماء.}{}

\Couplet{١٠٨}{«نعم، مات بان؛ وجاء الناصري وانتزع مقعده تحت الشمس؛}{عابد الإله اللغز، الذي واحده ثلاثة وثلاثته واحد.}{}

\Couplet{١٠٩}{«وعقيدته المحزنة في الخطيئة الموروثة صبّت على العالم سحرها الرمادي البارد؛}{وأظهرت في كل منظر قبرًا، وتحت القبر وهج الجحيم.}{}

\Couplet{١١٠}{«حتى يهبط شعر الحياة كله إلى نثر، وتبهت الرومانسية إلى واقع بليد؛}{ويشحب توهج الأرض بالفرح في الكآبة، ويهبط الإله من جديد إلى الإنسان.}{}

\Couplet{١١١}{«ثم جاء العربي النحيل، الملوث بالعرق، المستنزف لضرع الناقة؛}{المتخم بلحم سحليته الخضراء كالكراث، المكتسي بخرقته وبساطه القذرين.}{}

\Couplet{١١٢}{«حمل إلهه اللهَ الشرس عبر رماله، واندفع كانفجار الحمم}{إلى الممالك التي حكم فيها ملوك ما قبل آدم، والتي قام فيها العرش الكياني العظيم.}{\BurtonNote{حاشية برتون: الكيانيون من سلالة كورش، وهم أبطال المجوس القدماء.}}

\Couplet{١١٣}{«من يكترث اليوم للغناء عن كيومرث القديم، أو زال أو رستم؛}{وقد غمرتهم عاصفة القبائل التي سمّت سائق الجمل ملكًا؟}{}

\Couplet{١١٤}{«أين تاج كيخسرو، وصولجان أنوشروان؛}{وأين الكأس المقدسة لجمشيد الرفيع، وقاعة أفراسياب؟ أتستطيع أن تخبرني، أيها الإنسان؟}{}

\Couplet{١١٥}{«ذهبت، ذهبت إلى حيث يجب أن نذهب أنا وأنت، تحملها أجنحة الموت المذرّية؛}{ذلك الرعب الجاثم فوق الحياة، والذي يقترب مع كل نفَس.}{}

\Couplet{١١٦}{«لقد ملأ صيتهم الأقاليم السبعة؛ نهضوا وحكموا، وقاتلوا وسقطوا؛}{كما يعلو عبر البرية ثم يخفت رنين جرس الجمل».}{}
